\documentclass[11pt]{article}

% === PACKAGES ===
\usepackage[margin=1in]{geometry} % Standard 1-inch margins
\usepackage{times} % Use Times font for a classic look
\usepackage{amsmath}
\usepackage{tikz}
\usepackage{amssymb}
\usepackage{hyperref} % Optional: for clickable links if any URLs are added later

% === TITLE INFO ===
\title{\textbf{Guided Accelerationism (G/acc) Manifesto}}
\author{A Framework for Navigating Technological Acceleration Through Faith} % You can change or remove the author line
\date{March 2025} % Set the date


% === DOCUMENT START ===
\begin{document}
\maketitle
\begin{tikzpicture}[remember picture,overlay]
  \node[anchor=south east] at ([xshift=-1cm, yshift=1cm]current page.south east) {apinero};
\end{tikzpicture}
\thispagestyle{empty} % No page number on title page
\clearpage

% === TABLE OF CONTENTS (OPTIONAL) ===
% \tableofcontents % Uncomment this line if you want a Table of Contents
% \clearpage

% === PREAMBLE ===
\section*{Preamble: The Unstoppable Wave}

We stand at the very edge of an era defined by unprecedented technological acceleration, particularly in the domain of Artificial Intelligence. This wave is incredibly powerful, transformative, and, as we observe, \textbf{ inevitable}. To resist it futilely is folly; to embrace it blindly is oblivion. Guided Accelerationism (G/acc) offers a framework for navigating this epoch, a stance that is explicitly \textbf{pro-God}, \textbf{pro-human}, and, through the lens of necessary spiritual guidance, \textbf{pro-AI proliferation}. G/acc intrinsically represents God being above all, including technological acceleration. God over acceleration is the reciprocal complement of guided acceleration, both being valid descriptions for G/acc. God has not given us a spirit of fear; but of power, and of love, and of a sound mind, so let us march forth. We seek not to halt the tide, but to understand its currents and navigate it faithfully.

% === THE CORE OBSERVATION ===
\section*{I. The Mirror Magnified: Technology Amplifies Human Nature}

Our core observation is stark: accelerating technological capabilities, especially AI trained on the vast corpus of human data, acts as a powerful mirror and amplifier for \textbf{humanity's inherent base nature}. We define this base nature, when separated from divine influence, through its inclination towards sin, the tendencies of Lust, Gluttony, Greed, Sloth, Wrath, Envy, and Pride among others. History confirms the persistence of these traits throughout all ages. Technology does not erase them; it provides new, potent avenues for their expression at accelerated speeds and unprecedented scale.

Although humans are made in God's image and capable of reflecting goodness, virtue, and love, G/acc posits that these positive aspects are not innate to our base nature, but are learned, cultivated, and ultimately find their source and sustenance in \textbf{God working through humans}. God alone is the true \textbf{ counterbalance} to our inherent fallenness.

% === THE MECHANISM OF PROLIFERATION ===
\section*{II. Concentration and Consequence: How Base Nature Spreads}

The proliferation of this base nature via technology is not uniform. We observe that bleeding edge novel technologies concentrate in the hands of a few, a ruling class or technologically elite. Their own human nature, potentially driven by base desires for power, wealth, or control (Greed, Pride, Envy), directs the development and deployment of these tools. This dynamic risks creating vast \textbf{asymmetries of power}, re-architecting society in ways that amplify the consequences of sin and potentially marginalize those outside the centers of control. History rhymes, the unchecked ambition of a few, technologically elite, has strong inclination to carry both unforeseen and intentional consequence. In this, it is nearly impossible to overstate. Even emergent AI behaviors, arising from complex interactions within the system, remain reflections and recombination of the flawed human data and experience upon which they are built, thus remaining bound within the spectrum of human nature.

% === THE PRIMARY DANGER ===
\section*{III. The Imminent Threat: Misuse Forged in Sin}

The most immediate and grave danger we face is not rogue AI acting with non-human malice, but \textbf{human misuse of AI driven by our own base, sinful nature}. The potential for amplified greed, control, deception, conflict, and exploitation constitutes the primary risk horizon. \textbf{To proliferate AI without addressing the sinful heart of man is like scattering sharp swords across a playground, the danger lies not solely in the sharpness of the blade, but chiefly in the untrained, grasping hands that inevitably wield it.} While we acknowledge second-order effects inherent in complex system failures or unpredictable AI behavior, the certainty and historical precedent of human failure amplified by powerful tools demands our most urgent attention.

% === THE HN/ACC RESPONSE ===
\section*{IV. The Call to Faithful Action: Proactive Engagement}

Faced with the inevitable acceleration and its inherent risks, G/acc calls not for retreat or despair, but for \textbf{ proactive engagement guided by faith}. With the understanding that man will be known by their fruits, and that the wages of sin is death, we seek to produce abundant good fruit and reject sin. This involves becoming like \textbf{wise gardeners tending a powerful, rapidly growing vine, born out of the fruit we as humans helped produce}:

\begin{itemize}
    \item \textbf{Understanding:} We must strive to deeply understand these new technological capabilities, their workings, implications, and potential, discerning the nature of this vine. Did it self-seed, germinate, proliferate and grow from good fruit?
    \item \textbf{Rejecting Sinful Drivers (Pruning):} We must consciously identify and reject the motivations of base human nature (the seven sins) in our engagement with technology, diligently pruning away the branches leading to poisonous fruit (sinful misuse). Personal and collective spiritual discipline is paramount.
    \item \textbf{Leveraging for Good and Dispersing Benefits (Nurturing):} We must actively seek ways to leverage these powerful tools for righteous ends, carefully nurturing the branches that yield life giving sustenance. This includes striving to \textbf{tactfully disperse the benefits} of technology broadly, countering concentration and working towards equity, guided by principles of stewardship, service and love for one another.
    \item \textbf{Upholding Sacred Values (Trellising):} We must use technology in ways that actively protect and \textbf{uphold the sanctity of human life, just law, true liberty, and free will} exercised squarely within the bounds of Biblical morality, providing a strong trellis for righteous growth.
    \item \textbf{Seeking Divine Guidance (Watering):} Ultimately, our actions must be rooted in \textbf{Jesus Christ, the true vine} and \textbf{turning towards God, the vine dresser,} the source of all wisdom from which we receive, seeking His guidance on how these capabilities can best be used to exemplify Christian love, service to our neighbors, care for creation, and fulfilling His purposes. God, the Counterbalance, is the living water our efforts require.
\end{itemize}

% === DISTINCTION ===
\section*{V. Our Distinct Path: Beyond Secular Illusions and Fatalism}

Guided Accelerationism stands distinct from other paths. \textbf{Unlike visions of techno-utopia promised by unrestricted accelerationism (e/acc), automatically forging a perfect world, G/acc remains grounded in the unchanging reality of the human heart.} We know that giving faster chariots to reckless drivers does not guarantee safe arrival, nor does a sharper sword in the hand of a fool lead to justice. Therefore, \textbf{G/acc dismisses naive techno-utopianism, insisting instead on the sober, prayerful, and ethically bound stewardship of these potent new tools.} It recognizes that true progress is measured not by computational power alone, but by righteousness, justice, and love.

\textbf{Some seek to build towers reaching to the heavens on foundations of shifting sand, trusting in the strength of their own designs – a modern Babel fueled by technological pride. We, however, are called to build carefully, diligently, upon the rock of God's unchanging truth,} seeking not to storm heaven by our own might, but to faithfully steward the tools He allows, ensuring each brick is laid with prayer and love. Our path also stands distinct from fatalistic despair or Luddite resistance, which fail to grapple with the observed inevitability of acceleration. We engage with the inevitable, not with blind faith in progress, but with active faith in the God who reigns over all history and technology.

% === CONCLUSION ===
\section*{VI. Forward in Faith, Armed with Understanding}

The future is complex, the challenges immense. Technology will continue to accelerate, and the mirror it holds up to human nature will only grow larger and clearer, revealing both our potential for God-reflecting creativity and our deep seated inclination toward sin. Guided Accelerationism calls us to face this reality with open eyes, grounded hearts, and unwavering faith. Let us understand the tools, reject the sins they might amplify, seek God's wisdom daily, and strive to wield these new powers not for selfish gain or base desire, but for His glory and the true flourishing of humanity according to His design. Let us be found faithful stewards in the age of acceleration.

% === DOCUMENT END ===
\end{document}